\documentclass[../main.tex]{subfiles}
\begin{document}
%==========================================TIÊU ĐỀ TẬP========================================
\begin{table}[h]
  \begin{adjustwidth}{-1cm}{}
    \begin{tabular}{>{\centering\arraybackslash}p{6cm}|>{\raggedright\arraybackslash}p{12.5cm}}
      \multirow{5}{6cm}
      % Cột bên trái: ÉPISODE và số tập
      {\\[-40pt]
        \flaregothic\fontsize{20pt}{20pt}\selectfont \centering
        \color{eptype}ÉPISODE\color{black}\\
        \burbank\fontsize{72pt}{72pt}\selectfont
        \color{epnum}131
        \\[18pt]
        % \flaregothic\fontsize{14pt}{14pt}\selectfont
        % \color{epattr}BIENVENUE, \uline{\color{Banpire} Mel} !
        % \flaregothic\fontsize{18pt}{18pt}\selectfont
        % \color{epattr}ÉPISODE PERDU
        \color{black}
      }
       &
      % Dòng thứ nhất: tên tập tiếng Nhật
      {\fotskip\fontsize{18pt}{18pt}\selectfont \ruby{暑}{あつ}\ruby{苦}{くる}しい！！！
      }                                            \\[6pt]
      % Dòng thứ hai: tên tập tiếng Anh
       & {\cafeteria\fontsize{18pt}{18pt}\selectfont Don't You Dare be Sour!}                                           \\[4pt]
      % Dòng thứ ba: tên tập tiếng Việt
       & {\vnmsans\fontsize{18pt}{18pt}\selectfont Thám tử tư Matsuri}                                                  \\[8pt]
      % Dòng thứ tư: Nhân vật xuất hiện
       & {\sen{\fontsize{14pt}{14pt}\selectfont Track used: Mental Rave (\ruby{沸}{わ}き\ruby{立}{た}つ\ruby{祭}{まつり}/パーティー)}}
      \\[6pt]
      % Dòng cuối cùng: Thời gian diễn ra của tập (thường sẽ là ngày phát hành)
       & {\continuum\fontsize{16pt}{16pt}\selectfont Ngày phát hành: 14/11/2021}                                        \\[8pt]
       & (dựa theo bản dịch của \href{https://www.facebook.com/mamisubteam}{MamiSubs - Mami Fansub})                    \\[18pt]
    \end{tabular}
  \end{adjustwidth}
\end{table}
%=========================================TRUYỆN CHÍNH========================================
\color{black}

Một ngày trung tuần tháng Mười Một. Mặt trời lười biếng hơn mọi khi, lẩn khuất sau những đám mây xám dày đặc của mùa cuối thu. Không khí bên ngoài mang theo chút lạnh se sắt nhưng dễ chịu, một kiểu thời tiết dần sang đông khiến người ta chỉ muốn nằm cuộn trong chăn cả ngày. Thế nhưng ở văn phòng \hll, nơi cái sự ``êm đềm'' luôn bị đánh đổi bằng những tai nạn bất thình lình, thì khung cảnh ấy chỉ là vẻ ngoài đánh lừa.

Mọi thứ diễn ra có vẻ êm đẹp- à không, hiếm khi nào lại có được một ngày yên bình thực sự ở đây. Dù chẳng tới mức hỗn loạn như hồi họ dựng phim cho Mio-loli (một ác mộng hậu trường khó thể quên), vẫn có những tình huống dở khóc dở cười xảy ra bất chợt, như một định mệnh bất biến của hội chị em này.

Và tất nhiên, trong hoàn cảnh như thế, ai sẽ là người xắn tay áo giải quyết mọi chuyện?

``\textbf{Natsuiro Matsuri, thám tử tư kiêm tư vấn viên sẵn sàng phục vụ bạn đây!}'' tiếng reo lanh lảnh vang lên, đánh dấu màn tự xưng mới nhất trong vô vàn vai trò mà nàng cổ động viên đã thử qua.

\img{10-10a}

\dots Hoặc, ít nhất thì chính cô nghĩ như vậy.

Kuon vừa mới bước vào văn phòng và nghe thấy câu chào hoành tráng ấy, chỉ biết thở dài, như đã nghe những dạng câu này hàng trăm lần: ``Lại nữa à? Cái ngành `tư vấn viên' này hot vậy sao?''

Nàng Lễ hội ngạc nhiên khi nhìn thấy vẻ mặt chán chường của cậu: ``Sao anh nghe cứ như không mặn mà thế?''

Cậu Tổng thở dài, ngồi kể lại những trường hợp tương tự như vậy hồi trước: ``Anh từng thấy Shion-chan với Pekora-chan cũng thử nghề này rồi đấy. Một bên thì phán bừa lung tung, xong rồi còn sợ cả máu\footnote{Xem lại {\flaregothic ÉPISODE 108}, quyển số Tám.}, một bên thì lại đưa ra mấy cái lời khuyên làm hại người ta\footnote{Xem lại {\flaregothic ÉPISODE 111}, quyển số Chín.}\dots\ Và cái kết thì như em đoán rồi đấy, như cái đầu đinh luôn.''

Matsuri khựng lại một giây khi cậu dội thẳng gáo nước lạnh vào mong muốn của cô. nhưng rồi vẫn tươi tỉnh, quyết tâm làm ra trò. ``Thế nên em mới muốn thử! Em là Chúa Matsuri đấy nhá!''

``Chúa nỗi gì\dots'' Kuon chưng hửng chống hông, đảo mắt như nhìn trước kết quả. ``Đừng có quay về khóc lóc bảo `khách dỗi' khi thất bại đấy.''

``Rồi rồi\~{} Em sẽ cố gắng hết sức, cứ tin em đi!''

\ct

Vụ đầu tiên ngay lập tức được ghi nhận trong bếp văn phòng. Noel đang hí hửng chuẩn bị bữa trưa của mình, hộp mì yakisoba thịt bò yêu thích, nấu đúng theo công thức ``vừa đủ ăn vừa không bị Suisei chê nhạt.''

Nàng Hiệp sĩ vừa ngân nga vừa pha mì: ``Sắp có mì bò ăn rồi\~{} Mong chờ quá đi\~{}''

Kuon đứng ở bên cạnh, lên tiếng nhắc nhở: ``Cẩn thận không lại vụng tay vụng chân nữa nha. Hôm nọ trượt tay làm đổ cả nồi canh lên chân rồi đấy.''

``Anh cứ lo xa thế, em sẽ ổn thôi mà-''

*SOẠT*

Ngay khi cô chưa kịp dứt câu, trong một thoáng bất cẩn vì không dán kĩ nắp, toàn bộ mì cùng nước xúp chảy tong tỏng xuống bồn rửa, trôi theo dòng nước vào cống như một giấc mơ tan biến.

\img{10-10b}

Noel đứng chết lặng, mắt đờ đẫn như cá ươn, tay cầm cái nắp hộp trống rỗng. Chỉ thốt ra được một từ duy nhất: ``Ơ\dots''

Kuon lắc đầu thở dài, giọng nửa nghiêm nửa mỉa mai: ``Vừa mới nói xong. Thôi, thế này khỏi ăn luôn chứ sao nữa.''

Trong nỗi tuyệt vọng của cô nàng Thế hệ thứ Ba, Matsuri, như một vị cứu tinh tự phong, xuất hiện đúng lúc với vẻ mặt đầy hăng hái: ``Em thấy có ổn không?''

Cậu Tổng rướn mày: ``Em nhìn mặt em ấy thế kia mà còn hỏi nữa? Ít nhất cũng phải làm gì giúp người ta chứ?''

Noel tập tễnh cầm hộp mì rỗng không, đôi mắt rơm rớm, méo xẹo nói: ``M-Matsuri-senpai à\dots''

Không cần để cô nói thêm câu nào, nàng Lễ hội đột ngột chìa ra một hộp mì khác: ``Ăn cái này đi là em ổn liền! Mỳ siêu cay đặc biệt, phiên bản giới hạn cuối cùng đấy nhé!''

Noel nhìn hộp mì đỏ rực như vừa bước ra từ địa ngục, mặt càng xệch xuống hơn nữa, mắt nheo lại, khuôn mặt xanh xao hơn hẳn. ``Ờm\dots\ Em\dots\ Em không biết có nên\dots?''

\begin{figure}[H]
  \centering
  \begin{subfigure}[t]{0.45\textwidth}
    \centering
    \includegraphics[width=8cm]{../../../../Image/Book?/10-10c1.png}
  \end{subfigure}
  \quad
  \begin{subfigure}[t]{0.45\textwidth}
    \centering
    \includegraphics[width=8cm]{../../../../Image/Book?/10-10c2.png}
  \end{subfigure}
\end{figure}

``Anh nghĩ còn hơn là không có gì,'' Kuon nhún vai. ``Trừ khi em không ăn được cay.''

Tới đây, nàng Hiệp sĩ chỉ thẫn thờ đặt lại hộp mì xuống, quay lưng đi chậm rãi với nét mặt tiu nghỉu hiện rõ, để lại Matsuri đứng có cười tươi mà chẳng có ai hưởng ứng.

\ct

Một vụ việc khác xảy ra vào buổi sáng mù sương. Pekora đang ngáp ngắn ngáp dài trên đường đến văn phòng. Đi bên cạnh cô là Kuon, người đồng hành bất đắc dĩ trong những chuyến đi bộ sớm.

``Đêm qua lại stream khuya nữa hả?'' cậu Tổng nhướng mày, tỏ vẻ nghi ngại cho sức khỏe của cô.

``Đến tận hai giờ sáng mới ngủ đó, peko\dots'' nàng Thỏ đáp lại với một cái ngáp dài nữa.

Kuon liếc mắt nhìn phía trước, thấy một thứ gì đó nhô lên trên mặt đường: ``Này Peko-chan, ở phía trước có cái nút gì kìa-''

Chưa kịp cảnh báo, Pekora đã dẫm lên cái nút đó. Một tiếng *CẠCH* vang lên, một mảnh đất hình vuông mở ra như bẫy hố. Nàng Thỏ di chân phải tới trước, nơi mà đáng lẽ là đường đi bộ.

``\dots Ể?''

Trước khi cô có thể phản ứng, nàng Thỏ đã lọt tòm xuống cái hố sâu. Rất may là hố chỉ sâu chừng ba, bốn mét, được lót đất mềm, nhưng cô nàng vẫn đáp bằng\dots\ mông.

\begin{figure}[H]
  \centering
  \begin{subfigure}[t]{0.45\textwidth}
    \centering
    \includegraphics[width=8cm]{../../../../Image/Book?/10-10d1.png}
    \caption{Pekora vô tình giẫm vào nút công tắc\dots}
  \end{subfigure}
  \quad
  \begin{subfigure}[t]{0.45\textwidth}
    \centering
    \includegraphics[width=8cm]{../../../../Image/Book?/10-10d2.png}
    \caption{\dots khiến cánh cừa bẫy sập bật mở\dots}
  \end{subfigure}
  \quad
  \begin{subfigure}[t]{0.45\textwidth}
    \centering
    \includegraphics[width=8cm]{../../../../Image/Book?/10-10d3.png}
    \caption{\dots làm cho Pekora ngã dập mông.}
  \end{subfigure}
\end{figure}

Pekora xoa mông, rít lên một tiếng vì rát: ``Sống rồi-- mà ê mông quá peko! Sao lại có cái hố giữa đường thế này!?''

Kuon ở trên cao, ngó đầu vào đáp vọng: ``Anh cũng đâu biết. Thằng ranh nào cài bẫy chắc định chọc ai đấy.''

``Dù sao thì,'' cô hướng mặt lên trên, giọng có chút bối rối, ``'Tìm cách cứu em ra khỏi đây đi!'' Cô nàng nghiến răng đầy tức giận: ``Rồi em sẽ tóm cổ cái thằng ranh con đó cho mà xem!''

Một giọng nói hào hứng khác chen ngang, tông có phần nhấn nhá nổi bật hơn: ``Không cần đâu, có chị đây rồi!''

Không ai khác ngoài ``Thần Thám Tử'' (tự xưng) Matsuri, người đang thò đầu xuống cùng với cậu con trai. ``\textbf{Em thấy có ổn không!?}''

Cô nàng dưới dáy hố, trợn mắt quát lên: ``\textbf{Chị nghĩ em trông ổn chắc hả, peko!?}''

Câu hỏi của nàng Lễ hội còn khiến cho cậu Tổng không nhịn được mà nhịn cười: ``Em hỏi hơi ngu đấy\dots''

Dù thế, Matsuri không hề nao núng. Cô lấy ra một củ cà rốt to tướng, giơ lên như thể đó là bảo vật. ``Nhận lấy đi! Rồi em sẽ thấy khá hơn!''

Ngay lập tức, mắt nàng Thỏ sáng như đèn pha. Cô ngồi bật dậy, vươn tay với lấy, nhưng nàng Lễ hội lại buộc củ cà rốt vào một sợi dây, rồi bắt đầu thả xuống như đang đi câu cá.

Mỗi lần Pekora nhảy lên chụp lấy, củ cà rốt lại bị giật lên.

\img{10-10e}

Lần một.

Lần hai.

Lần ba.

Cho tới lần thứ tư, biết rằng mình đã bị chơi xỏ, nàng Thỏ tức giận, gân hiện rõ trên nét mặt hét lớn: ``\textbf{Oi!! Đùa đủ rồi đấy, Matsuri!!}''

Matsuri thì cười nghiêng ngả với phản ứng cay cú của đàn em, còn Kuon đứng bên cạnh lắc đầu chịu thua, chống hông hỏi: ``Thế rốt cuộc em định giúp hay trêu người ta đấy?''

Cô nàng không đáp, mà xách cần câu chạy đi mất, vừa chạy vừa cười khà khà.

``Ơ này!? Không định giúp thật à!?'' cậu đuổi theo cô, vô tình bỏ quên lại nàng Thỏ dưới hố sâu.

``\textbf{Hai người quay lại mau!! Kuon-senpai!! Đừng bỏ em ở đây một mình mà, peko!!}''

\ct

Một vụ tiếp theo diễn ra tại sân sau của văn phòng lúc hoàng hôn. Suisei đeo mặt nạ Jason Voorhees từ Friday the 13th và chiếc cưa máy đang chuẩn bị cho một cảnh quay.

Đối diện cô là Shion, nay hóa thân thành ``thám tử lừng danh'', tay cầm cặp hồ sơ, miệng cười đắc thắng: ``Bằng chứng ngoại phạm của ngươi đã bị phá huỷ rồi! Đầu hàng đi, Suisei!''

\begin{figure}[H]
  \centering
  \begin{subfigure}[t]{0.45\textwidth}
    \centering
    \includegraphics[width=8cm]{../../../../Image/Book?/10-10f1.png}
  \end{subfigure}
  \quad
  \begin{subfigure}[t]{0.45\textwidth}
    \centering
    \includegraphics[width=8cm]{../../../../Image/Book?/10-10f2.png}
  \end{subfigure}
\end{figure}

Nàng Ca sĩ hơi khựng lại. Có vẻ như cô không hoàn toàn hiểu Shion đang nói đến cái gì, bởi ``kế hoạch'' của cô, nếu có, cũng chỉ là dọa Ma-chama (Haato).

``Hả!? Ngươi đang nói cái gì vậy!?'' Cô vung cái máy cưa lên như định truy đuổi tên ``thám tử Tỏi'', thì\dots

*PHÀNH PHẠCH* *PHÀNH PHẠCH* *PHÀNH PHẠCH*

Tiếng máy bay trực thăng rền vang từ trên cao. Một sợi dây thang rơi xuống, kéo theo Matsuri đang tươi cười đáp xuống như siêu anh hùng.

``\textbf{Em thấy có ổn không!?}''

Tiếng quạt quay ồn ào tới độ đối phương không thể nghe rõ, ghé sát tai hơn: ``Hả? Nói gì cơ?''

Từ trên máy bay, người phi công lái trực thăng (lại là Kuon) đáp lại từ bộ đàm: ``Bỏ cái mặt nạ ra thì nghe rõ hơn đó!''

``À, ừ\dots\ Đợi chút.''

Cô tháo mặt nạ Jason, rồi quay lại nhìn Matsuri. ``Em vừa nói gì thế?''

Nàng Lễ hội đưa tay về phía cô, tươi cười nói: ``Trốn cùng em đi! Rồi sẽ ổn hết!''

Cậu Tổng hướng dẫn qua bộ đàm: ``Cứ bám lấy thang dây là được. Phần còn lại anh sẽ lo liệu.''

Không chần chừ, Suisei nắm lấy dây thang, leo lên máy bay. Matsuri theo sau. Trực thăng cất cánh, để lại Shion, người vừa bị ``bắt hụt tội phạm'' tức giận. Cô vừa chạy theo, vừa giơ nắm đấm như nguyền rủa, vừa hét lên: ``\textbf{Ta bảo mấy người đứng lại đó mà!! Đồ nhát cáy!!}''

\img{10-10g}

Kuon lẩm bẩm bằng giọng mệt mỏi: ``\textit{Sao ình cảm tường đang tiếp tay cho tội phạm thế nhỉ?}''

Ở trên thang dây, nàng Lễ hội vẫy tay tới người đang chạy ở dưới, không khác gì đang chọc tức vị ``thám tử'' ấy: ``Em đang giúp chị ấy mà!''

Nàng Ca sĩ giơ nắm đấm lên cao, vui mừng như trúng số: ``\textbf{Tuyệt quá! Cảm ơn hai người nha!}''

Cuộc rượt đuổi giữa Shion và chiếc trực thăng cứ thế tiếp diễn, cho tới khi máy bay lên cao tới độ cô nàng không thể đuổi được nữa. Tiếc nuối khi miếng mồi ngay trước mắt đã bị vuột mất, cô ngửa cổ hét lớn: ``\textbf{Khôôôôôônggggg!!!!}''

\ct

Một vụ kỳ lạ khác xảy ra vào xế chiều hôm đó, khi Matsuri vẫn hăng hái trong vai trò ``tư vấn viên/thám tử'' của mình lại cùng Kuon lảng vảng quanh văn phòng, với mục đích (tự phong) là ``giúp đỡ những ai cần đến bàn tay Thần Thánh  của mình.''

Nhưng thực chất ra thì chỉ có mình cô mới hớn hở như thế.

``Anh tưởng em nói giúp người, chứ có thấy ai nhờ em đâu?''

``Chưa nhờ không có nghĩa là không cần giúp!''

Câu trả lời đầy nhiệt huyết của cô làm cho cậu Tổng thở dài, khoanh tay lại nhìn nàng Cổ động như một người giám sát mệt mỏi. Còn cô thì bắt đầu tìm khắp nơi trong văn phòng, lật tung cả đệm ghế, bới cả dưới bàn, thậm chí là dí sát tai vào bức tường.

``Tìm người thôi chứ có phải đào kho báu đâu\dots'' Kuon toát mồ hôi hột.

Bất chợt, Matsuri dừng lại trước ghế sofa lớn.

Ở đó, ba thành viên là Mel, Roboco và Botan đang bất động. Roboco và Mel ngả nghiêng ngủ gật trên ghế, còn Botan thì nằm úp mặt giữa phòng như cái đống lông xám trắng đổ ụp. Nhìn bề ngoài thì chẳng có gì bất thường, có thể họ đơn giản chỉ là ngủ trưa quá giấc.

\img{10-10h}

Matsuri tiến lại gần, cúi xuống ``quan sát hiện trường.'' Cô sờ trán từng người, kiểm tra hơi thở (vẫn ổn), rồi đứng thẳng dậy kết luận: ``Không có gì khả nghi cả\dots\ hoặc là quá yên tĩnh để khả nghi.''

``Anh nghĩ là chúng ta kết thúc trò mèo này ở đây thôi-''

Nhưng ngay khi cậu còn chưa nói hết câu, một ánh lóe nhẹ phía sau đầu Roboco làm cô chú ý. Từ cái ahoge (chỏm tóc ngố) đặc trưng của nàng Người máy, một đường ánh sáng mờ mờ, mảnh như chỉ, bắt đầu uốn lượn trong không khí.

``\dots Cái gì thế này? Hiện tượng siêu nhiên à?''

``Cả Mel-chan và Botan-chan cũng có một vệt tương tự từ chỏm tóc kìa. Hơi\dots\ kỳ quặc đấy.''

\img{10-10i}

Cả hai nhìn theo cái đường ánh sáng lấp lánh, như một luồng khí ma quái nhưng\dots\ không hề đe dọa. Nó nhẹ nhàng dẫn ra khỏi văn phòng, như thể đang chỉ đường cho họ.

Và tất nhiên, như những logic trong các phim trinh thám mà ta đã từng biết, không dễ gì mà Matsuri có thể bỏ qua chúng.

``Đi theo!''

``\dots Biết ngay là sẽ thế mà.''

Cô rút kính lúp ra (dù chẳng giúp ích gì), bắt đầu dò theo đường sáng như thể đang chơi trò ``Đuổi hình bắt bóng.'' Kuon lững thững theo sau, tay chắp sau lưng như một giám thị cho thực tập sinh.

Họ đã đi qua\dots

\begin{figure}[H]
  \centering
  \begin{subfigure}[t]{0.45\textwidth}
    \centering
    \animategraphics[autoplay,loop,width=8cm]{12}{../../../../Image/Book?/10-10j/10-10j.}{1}{146}
    \caption{\dots hết văn phòng\dots}
  \end{subfigure}
  \quad
  \begin{subfigure}[t]{0.45\textwidth}
    \centering
    \animategraphics[autoplay,loop,width=8cm]{12}{../../../../Image/Book?/10-10j/10-10j.}{200}{270}
    \caption{\dots theo từng con hẻm nhỏ trên phố\dots}
  \end{subfigure}
  \quad
  \begin{subfigure}[t]{0.45\textwidth}
    \centering
    \animategraphics[autoplay,loop,width=8cm]{12}{../../../../Image/Book?/10-10j/10-10j.}{330}{470}
    \caption{\dots băng qua khu rừng sinh thái rộng lớn\dots}
  \end{subfigure}
\end{figure}

\dots cuối cùng lạc luôn vào một bãi thảo nguyên hoang vu. Mây đen bắt đầu kéo tới, một lớp sương mù mỏng phủ quanh chân cỏ, khiến khung cảnh trông vừa kỳ ảo vừa hơi lãng đãng.

Ở giữa thảo nguyên ấy là một cây cổ thụ khổng lồ, xù xì và âm u như thể tồn tại từ thời tiền sử.

\img{10-10k}

``Cái vệt sáng đó kết thúc ở đây\dots'' Kuon ngước nhìn lên cái cây, trầm ngâm suy nghĩ một lúc lâu. ``Anh không chắc liệu nó là manh mối hay không nữa.''

Matsuri khoanh tay lại, xâu chuỗi lại những gì vừa xảy ta từ lúc họ thấy ba người họ ở văn phòng tới đây. ``Có khi nào cái cây này là đầu mối cuối cùng không?''

``Ừm. Nếu cái cây đó biết nói thì chắc đúng. Nhưng mà anh thấy--''

Đúng lúc cậu đang nghĩ ra một câu phản biện, cái cây cổ thụ rung nhẹ, và bất ngờ phát ra một giọng nói trầm mặc vang vọng: ``\emph{Có chứ.}''


Cả hai đứng sững sau giọng nói trầm, vang, nhưng gây ấn tượng đó. Không phải ấn tượng vì một sự cứng cáp, mà bởi sự xuất hiện bất ngờ của nó.

Gió thổi nhẹ. Mây trôi chậm. Đất không rung, nhưng mặt Kuon thì giật: ``Nó nói chuyện thật à? Chết tiệt, lại thêm cái vụ huyễn hoặc gì nữa đây?''

Trái với một cậu Tổng khẽ hoảng loạn, nàng Lễ hội dường như vẫn khá bình tĩnh. Cô tiếp cậ tới cái cây già, cúi đầu nhẹ trước nó như một vật linh thiêng và hỏi: ``Ta chỉ muốn hỏi\dots\ Ngươi có ỔI (ổn) không?''

\dots

\dots

Im lặng.

Tiếng gió xào xạc thổi thoáng qua, những tiếng lá rơi cọ xát với nhau, đi cùng với bầu không khí âm u của buổi chiều tối cuối thu, càng làm cho mọi thứ huyền ảo hơn.

Và rồi, như để đáp lại sự kiên nhẫn tới đáng sơ của thiên nhiên, cái cây đáp lại một lần nữa cũng  bằng cái giọng trầm kia: ``Có chứ.''

Kuon nhìn sang Matsuri, mặt đơ như tượng đá. Đầu óc cậu quay cuồng, cố gắng nghĩ tới vắt óc để hiểu được thông điệp mà cái cây đó đã gửi đến.

Nhưng cho đến cuối cùng, hóa ra đó chỉ là\dots\ một trò lừa đảo.

``\dots Đợi đã. Vậy nãy giờ là tất cả chỉ là cái trò chơi chữ ``Ổi/Ổn'' thôi sao!? Cái đường đó\dots\ chỉ dẫn tới một cây biết đáp đúng vần điệu thôi ư!?''

Nàng Lễ hội nở một nụ cười gian, như thể vừa chơi khăm cấp trên thành công: ``Chuẩn luôn\~{}''

Cậu Tổng ngã chổng vó xuống đất, than vãn thấu cả trời xanh: ``Vậy mà anh cứ tưởng cái gì ghê gớm chứ\dots''

Đến cuối cùng, cái cây cổ thụ đó hoàn toàn không có mối quan hệ gì với bộ ba đang ngủ ở văn phòng cả!

Một vụ án nữa kết thúc với càng nhiều dấu hỏi đổ lên đầu cô nàng ``Thám Tử Tư kiêm Tư Vấn Viên'' này hơn, nhưng có lẽ chúng ta sẽ không bao giờ biết được cô nàng đang nghĩ gì trong đầu.
%=========================================TRUYỆN CHÍNH========================================
\omk

Sau vài ngày liên tiếp đóng vai tư vấn viên kiêm thám tử tự phong, Matsuri vẫn không hề nản chí, mặc cho những lời ca thán từ cac thành viên trong công ty. Trái lại, cô còn tìm đến Kuon với gương mặt rạng rỡ, mong đợi một lời khen: ``Anh thấy em làm một thám tử như thế nào, Kuon-senpai?''

Cậu Tổng im lặng vài giây, rồi trả lời một cách chung chung, không rõ là đang khen hay chê: ``Đỡ hơn Shion-chan, nhưng không thể gọi là tốt được.

``Ể?'' nàng Lễ hội nghiêng đầu. ``Ý anh là sao chứ?''

Cậu thở dài một hơi nữa, rồi đáp lại một cách gọn lỏn: ``Nói trắng phớ ra là không đạt. Có giúp, nhưng không đáng kể.''

Cô nàng giật mình, mắt tròn xoe khi bị chính giám thị đánh trượt tại chỗ: ``Ể\dots\ Nhưng mà em tưởng em đã làm tốt mà? Em có giải quyết vụ của Noel-chan đấy thôi.''

``Ừ thì đúng là có giải quyết, nhưng bằng cách dúi cho cô ấy một hộp mì siêu cay mà người ta còn đang khóc vì cái cũ rơi mất thì không gọi là tư vấn,'' Kuon phản biện lại bằng giọng như một giáo viên nghiêm khắc. ``Cái đó có thể quy vào dạng đầu độc người ta đấy.''

``C-Còn vụ của Pekora-chan thì sao?''

``Em bỏ mặc em ấy dưới hố, đã thế còn chọc tức người ta bằng cà rốt. Em làm còn chẳng đến nơi đến chốn kìa.''

Từ những lời đanh thép của cậu, Matsuri dần chùn bước, chấp nhận (một cách chóng vánh) với kết quả này. Cô gãi đầu gãi tai, nhớ lại về hai vụ mà cô đã xử trước đó: ``Thế rốt cuộc\dots\ Noel-chan và Pekora-chan giờ sao rồi nhỉ\dots''

Cậu Tổng tặc lưỡi một cái, rồi nói với giọng thất vọng: ``Dám chắc là sẽ giận em, nhất là Peko-chan đấy.''

``\dots Mà, anh nói là tốt hơn Shion-chan là ở khoản nào vậy?''

``Cứ nhìn em ấy xử lý vụ Peko-chan là hiể-u'' vừa nói đến đó, Kuon khựng lại như chợt nhớ điều gì, rồi đập tay lên trán: ``Mà nhắc đến cà rốt thì anh còn quên mất chưa lôi \emph{em ấy} ra khỏi hố đấy\dots''

\ct

Ở một góc phố vắng, vài ngày sau sự cố nút bấm đó.

Pekora vẫn đang bị kẹt trong cái hố bẫy suốt nhiều ngày, ánh mắt lờ đờ, tóc tai rối bù. Cô ngồi bệt dưới đáy, rầu rĩ lầm bầm: ``Kuon-senpai chạy đi đâu rồi không biết\dots\ Định bỏ mình ở đây cho chuột ăn à\dots''

Bất chợt, một cái đầu ló lên trên miệng hố. Đó là vị ``thám tử tư'' vừa vụt mất con mồi Suisei. Tay cầm tẩu thuốc, vẻ mặt đầy tự đắc, nàng Phù thủy nghiêng đầu hỏi: ``Có gì khiến em cảm thấy bất ổn à, nghi phạm?''

Nàng Thỏ ngước lên, giận tới tím mặt: ``\textbf{Tôi đang kẹt dưới HỐ đấy! Nhìn mặt tôi thế này mà chị còn hỏi à, peko!?}''

\img{10-10l}

Shion gật gù, như thể đã ghi lại lời khai quan trọng. ``Tình trạng cảm xúc không ổn định. Nghi phạm có dấu hiệu hoảng loạn. Khả năng che giấu sự thật cao.''

``\textbf{Chị có định cứu tôi khỏi đây không thì bảo!?}'' Pekora hét lên, lửa giận gần như bốc khói từ đôi tai thỏ đang quăn lại.

``Nếu như em thừa nhận tội thì ta có thể khoan hồng,'' Shion vẫn giữ nguyên nét mặt của một người bình thản, dù hoàn toàn không hợp tông với bối cảnh hiện tại.

``\textbf{Tội tình gì chứ, peko!? Tôi có làm gì sai trái đâu chứ!!}''

``Thái độ của cô càng lúc càng đáng nghi hơn đấy,'' nàng Phù thủy ghi chú thêm, nghiêm trọng hóa mọi thứ.

Sau vài lời qua tiếng lại khiến Pekora gần như phát điên, cô bắt đầu xả giận, đúng nghĩa đen:
``\textbf{Nếu tôi thoát ra khỏi đây được, tôi thề là tôi sẽ vặt sạch từng sợi tóc nhuộm tím của chị, peko!!}''

May mắn thay, vào thời điểm cô sắp bốc hỏa thành tro, Kuon cuối cùng cũng xuất hiện, mang theo một sợi dây thừng. Cậu buộc chặt dây rồi thả xuống, nhẹ nhàng ra hiệu: ``Bám vào đi, anh kéo em lên. Xin lỗi vì đến muộn.

Nhìn thấy chiếc phao cứu sinh - dù đã đến muộn - nàng Thỏ gần như phát khóc vì mừng. Sau khi được kéo lên khỏi hố, cô nằm phịch ra đất, thở dốc như người mới bơi từ đáy đại dương lên bờ.

Tuy nhiên, thay vì cảm ơn, cô lập tức bật dậy, chỉ tay mắng Kuon xối xả - cũng dễ hiểu tại sao: ``\textbf{Anh đi đâu mà để em ở dưới đó suốt mấy ngày vậy!!??}''

Cậu Tổng đồ mồ hôi, gượng cười: ``Cũng tại anh mải xử lý công việc của Mat-chan quá. Nhưng mà ít ra anh còn quay lại cứu mà, đúng chứ?''

Cu liếc về phía Shion đang chỉnh lại tẩu thuốc, và Matsuri vừa mới tới, đang huýt sáo như thể chưa biết chuyện gì xảy ra.

``\dots Tạm tha cho anh lần này, peko. Còn bây giờ\dots''

Không mất quá nhiều thời gian để Pekora quay sang mắng hai người tiền bối của mình như chưa bao giờ được mắng. Giọng cô to như sấm, lời lẽ sắc như dao bào, khiến cả hai người chỉ có thể vừa cười chát vừa lầm lũi nhận lỗi từ cô nàng ``tội phạm chiến tranh'' này.

Kể từ sau sự kiện đó, nàng Thỏ bắt đầu có ác cảm ngầm với tất cả những ai tự nhận mình là ``thám tử nghiệp dư'' trong \hll.

\dots Ngoại trừ người thực sự đã giúp cô kéo ra khỏi hố mà không gây thêm rắc rối nào khác\dots\ Chắc các bạn cũng đoán được đó là ai rồi đấy.

\end{document}